% from main.tex
%%%%%%%%%%  scan idx 573  |  folio 560  %%%%%%%%%%
\section*{7. Appendice : Spécialisation en théorie des intersections}
\begin{center}
par A. Grothendieck
\end{center}

\textbf{Lemme 7.1.} Soit \(i:Y\to X\) une immersion fermée régulière (VII 1.4), de codimension bornée, de sorte que \(i\) est un morphisme parfait (VII 1.9), et qu'on peut donc définir le morphisme (IV 2.12)
\[
 i^*:K_\bullet(X)\longrightarrow K_\bullet(Y).
\]
Soit d'autre part \(N\) le faisceau conormal pour \(i\), qui est un faisceau localement libre de type fini sur \(Y\) (VII 1.6), de sorte qu'on peut former
\[
 \lambda_{-1}(N)\in K^\circ(Y).
\]
Considérons aussi l'nomomorphisme image directe (IV 2.11)
\[
 i_*:K_\bullet(Y)\longrightarrow K_\bullet(X).
\]
Ceci posé, on a la formule
\[
 i^*i_*(x)=x\lambda_{-1}(N)\qquad\text{pour } x\in K_\bullet(Y),
\]
où le produit du deuxième membre est relatif à la structure de \(K^\circ(Y)\)-module de \(K_\bullet(Y)\) définie dans (IV 2.10).

La démonstration est, essentiellement sans changement, la même que pour la même formule, écrite en travaillant avec \(K^\circ(X)\), \(K^\circ(Y)\), établie dans (VII 2.7).

\textbf{Corollaire 7.2.} Soient \(f:X\to S\) un morphisme plat, avec \(S\) un trait i.e. le spectre d'un anneau de valuation discrète \(A\), \(s\) le point fermé de \(S\), \(t\) son point générique, \(i:X_s\to X\) l'inclusion canonique. Alors \(i\) est de tor-dimension finie donc (IV 2.12) on peut définir
\[
 i^*:K_\bullet(X)\longrightarrow K_\bullet(X_s),
\]
